\documentclass[10pt]{article}
\usepackage[margin=1in]{geometry}
\usepackage[T1]{fontenc}
\usepackage{lmodern,microtype,amsmath,amssymb,amsthm}
\theoremstyle{definition}
\newtheorem{assumption}{Assumption}
\begin{document}
\section*{Terminal Clipping for Convex Minimization}
We consider minimizing a measurable convex function $f:\mathbb{R}^d\to\mathbb{R}$ over a closed convex set $D$, with a given $x^0\in D$ within distance $R\geq0$ of a minimizer $x^\star$. The oracle returns stochastic subgradient estimates $g(x,\xi)$ from independent samples $\xi$ with law $\nu$, or exact projections onto $D$. Their means are subgradients on $D$, and $\mathbb{E}_{\xi\sim\nu}\|g(x,\xi)\|^2\leq B^2\|x-x^0\|^2+G^2$, where $B,G>0$. Projected clipped updates must return a measurable terminal iterate $\bar{x}\in D$ satisfying $f(\bar{x})-f(x^\star)\leq\varepsilon$ with probability at least $1-\delta$, where $\varepsilon>0$ and $\delta\in(0,1)$.

The optimization problem is
\[
\min_{x\in D} f(x).
\]
Here $d$ is a nonnegative integer, and $\langle\cdot,\cdot\rangle$ and $\|\cdot\|$ denote the Euclidean inner product and norm. The algorithm receives $d,x^0,B,G,\delta,R,\varepsilon$ and accesses the objective and constraint set through the oracle.

Our assumptions:
\begin{assumption}[Convexity and initial distance]
The set $D\subseteq\mathbb{R}^d$ is closed and convex, $x^0\in D$, and $f$ is measurable on $\mathbb{R}^d$. For all $x,y\in D$ and $t\in[0,1]$,
\[
f((1-t)x+ty)\leq(1-t)f(x)+tf(y).
\]
There exists $x^\star\in D$ such that
\[
\begin{aligned}
f(x^\star)&\leq f(x) &&\text{for every }x\in D,\\
\|x^0-x^\star\|&\leq R.
\end{aligned}
\]
\end{assumption}

\begin{assumption}[Stochastic subgradients]
Let $\nu$ be a probability measure on a measurable sample space $\Omega$. The map $g:\mathbb{R}^d\times\Omega\to\mathbb{R}^d$ is jointly measurable. For every $x\in D$, $g(x,\cdot)$ is Bochner integrable and $\|g(x,\cdot)\|^2$ is integrable under $\nu$. Its mean satisfies the subgradient inequality
\[
f(x)-f(y)\leq
\left\langle\mathbb{E}_{\xi\sim\nu}[g(x,\xi)],x-y\right\rangle,
\qquad x,y\in D,
\]
and its second moment satisfies
\[
\mathbb{E}_{\xi\sim\nu}\bigl[\|g(x,\xi)\|^2\bigr]
\leq B^2\|x-x^0\|^2+G^2,
\qquad x\in D.
\]
The oracle samples form a sequence $(\xi_k)_{k\geq0}$ of measurable, mutually independent random variables, each with law $\nu$. This sequence is independent of the algorithm's randomness. At call $k$, a subgradient query at $x$ returns $g(x,\xi_k)$.
\end{assumption}

\begin{assumption}[Exact projection]
A projection query at any $z\in\mathbb{R}^d$ returns $\Pi_D(z)$, where $\Pi_D:\mathbb{R}^d\to\mathbb{R}^d$ is measurable and satisfies
\[
\begin{aligned}
\Pi_D(z)&\in D,\\
\langle z-\Pi_D(z),u-\Pi_D(z)\rangle&\leq0,
\qquad u\in D.
\end{aligned}
\]
\end{assumption}

Our formulation is algorithm-specific: we analyze terminal clipping with projected stochastic subgradient updates. For a nonnegative threshold $C$, define
\[
\operatorname{clip}_{C}(v)=
\begin{cases}
v,&\|v\|\leq C,\\
\dfrac{C}{\|v\|}v,&\|v\|>C.
\end{cases}
\]
The method starts at the given $x^0$, with initial output $x^0$ and first query a subgradient query at $x^0$. It alternates subgradient and projection queries. Before receiving the subgradient reply at iteration $t$, it selects a step size $\eta_t\geq0$ and threshold $C_t\geq0$ using its available history and internal randomness. The updates are
\[
\begin{aligned}
g_t&=g(x^t,\xi_{2t}),\\
z^t&=x^t-\eta_t\operatorname{clip}_{C_t}(g_t),\\
x^{t+1}&=\Pi_D(z^t).
\end{aligned}
\]
The output remains $x^t$ after the subgradient reply and becomes $x^{t+1}$ after the projection reply; the next query is a subgradient query at $x^{t+1}$. Both query types count toward a single total oracle-call budget. Thus each completed update uses two calls. With a budget of $N\geq0$ calls, the terminal output is $\bar{x}=x^{\lfloor N/2\rfloor}$; an unmatched final subgradient call leaves it unchanged.

The required output is measurable and belongs to $D$ for every realization. Its accuracy requirement is
\[
\mathbb{P}_{P}\!\left(f(\bar{x})-f(x^\star)>\varepsilon\right)
\leq\delta,
\]
where $P$ is the joint law of the oracle samples and the independent algorithm randomness.
\end{document}
